\section*{Chapter \ref{ch:g12:internal-heat} --- The Earth's Internal Heat}

\begin{solution}{exo:g12:internal-heat:1}
(a) the \omterm{def:g12:internal-heat:heat-flow}{heat flow}; (b) the \omterm{def:g12:internal-heat:heat-production}{radiogenic heat production}; (c) \omterm{def:g12:internal-heat:secular-cooling}{secular cooling};
(d) \omterm{def:g12:internal-heat:convection}{thermal convection}; (e) \omterm{def:g12:internal-heat:geothermal}{geothermal energy}.
\end{solution}

\begin{solution}{exo:g12:internal-heat:2}
(1) \qty{0.0916}{W/m^2}. (2) \qty{4.67e13}{W}. (3) \qty{0.029}{K/m}: a
difference of one degree Celsius is a difference of one kelvin.
\end{solution}

\begin{solution}{exo:g12:internal-heat:3}
Thorium-232, uranium-238, potassium-40 and uranium-235. KamLAND sees the
geoneutrinos of thorium-232 and uranium-238 only.
\end{solution}

\begin{solution}{exo:g12:internal-heat:4}
$490/\sqrt{1} = \qty{490}{mW/m^2}$; $490/2 = \qty{245}{mW/m^2}$;
$490/5 = \qty{98}{mW/m^2}$; $490/7 = \qty{70}{mW/m^2}$, about the mean of the
continents.
\end{solution}

\begin{solution}{exo:g12:internal-heat:5}
Near the surface, the temperature of the ground follows the days and the
seasons; only below a few metres does it stay at the mean temperature of the
year, and only deeper still does it show the heat coming from below.
Drilling stirs the rock and the water in the hole and changes their
temperature, so one waits until the hole has come back to the temperature of
the rock.
\end{solution}

\begin{solution}{exo:g12:internal-heat:6}
(1) $(265 - 10)/9.1 \approx \qty{28.0}{\celsius}$ per kilometre, that is
\qty{0.0280}{K/m}. (2) $q = 2.72 \times 0.0280 \approx \qty{0.0762}{W/m^2}$,
about \qty{76}{mW/m^2}. (3) A little above the mean of the continents,
\qty{70.9}{mW/m^2}, by about \qty{7}{\percent}: a typical continental value.
\end{solution}

\begin{solution}{exo:g12:internal-heat:7}
Continents: $0.0709 \times \num{2.073e14} \approx \qty{1.47e13}{W} =
\qty{14.7}{TW}$. Oceans: $0.1054 \times \num{3.028e14} \approx
\qty{3.19e13}{W} = \qty{31.9}{TW}$. Total about \qty{46.6}{TW}, of which
$31.9/46.6 \approx \qty{68}{\percent}$, two thirds, leaves through the ocean
floor.
\end{solution}

\begin{solution}{exo:g12:internal-heat:8}
(1) $\num{1.88e-6} \times 8000 \approx \qty{0.0150}{W/m^2} = \qty{15.0}{mW/m^2}$.
(2) $15.0/70.9 \approx \qty{21}{\percent}$, a fifth of the \omterm{def:g12:internal-heat:heat-flow}{heat flow}.
\end{solution}

\begin{solution}{exo:g12:internal-heat:9}
$\sqrt{t} = 490/70.9 \approx 6.91$, so $t \approx 47.8$ million years: ocean
floor older than about 48 million years loses less heat than the mean
continent.
\end{solution}

\begin{solution}{exo:g12:internal-heat:10}
Uranium-235: $2^{1.4/0.704} = 2^{1.99} \approx 4$ times more. Potassium-40:
$2^{2.5/1.248} = 2^{2.00} \approx 4$ times more. Both are two half-lives back.
\end{solution}

\begin{solution}{exo:g12:internal-heat:11}
(1) $10^2/10^{-6} = \qty{1e8}{s}$, about 3 years. (2) $(10^5)^2/10^{-6} =
\qty{1e16}{s}$, about 320 million years. (3) $(\num{2.891e6})^2/10^{-6}
\approx \qty{8.4e18}{s}$, about 265 billion years. The time grows as the
square of the distance: ten times thicker, a hundred times longer.
\end{solution}

\begin{solution}{exo:g12:internal-heat:12}
Multiply each by $2^{3/T}$: thorium-232, $10.4 \times 2^{3/14.0} \approx
\qty{12.1}{TW}$; uranium-238, $9.6 \times 2^{3/4.468} \approx \qty{15.3}{TW}$;
potassium-40, $4.0 \times 2^{3/1.248} \approx \qty{21.2}{TW}$; uranium-235,
$0.42 \times 2^{3/0.704} \approx \qty{8.1}{TW}$. Total about \qty{56.6}{TW},
more than twice today's \qty{24.4}{TW}. The two short-lived isotopes give
$(21.2 + 8.1)/56.6 \approx \qty{52}{\percent}$ of it, against
$4.42/24.4 \approx \qty{18}{\percent}$ today.
\end{solution}

\begin{solution}{exo:g12:internal-heat:13}
The power is the mass per second times the specific heat times the
warming: $\num{3e8} = 1640 \times 4180 \times \Delta T$, so $\Delta T =
\num{3e8}/(1640 \times 4180) \approx \qty{44}{K}$. The average continent loses
\qty{0.0709}{W} per square metre, so \qty{300}{MW} is the natural heat loss
of $\num{3e8}/0.0709 \approx \qty{4.2e9}{m^2}$, about \qty{4200}{km^2}:
the plant gathers heat that the Earth lets out naturally over thousands of
square kilometres.
\end{solution}

\begin{solution}{exo:g12:internal-heat:14}
(1) 100 million years is $\num{2.6e11}/\num{1e8} = 2600$ times
shorter than 260 billion years; 200 million years 1300 times shorter.
(2) Convection carries the heat with the moving rock: the time to cross the
mantle is its thickness divided by the speed, $\qty{2891}{km}$ at a few
centimetres a year, about a hundred million years. Conduction has no such
carrier: its time grows as the square of the distance, and over thousands of
kilometres it becomes enormous.
\end{solution}

\begin{solution}{exo:g12:internal-heat:15}
(1) $10/22.3 \approx \qty{45}{\percent}$; the rest is the cooling of the
mantle itself, which loses its own stored heat. (2) The boundary lies at radius $6371 - 2891 =
\qty{3480}{km}$; its area is $4\pi (\num{3.48e6})^2 \approx
\qty{1.52e14}{m^2}$, and $\num{1e13}/\num{1.52e14} \approx
\qty{0.066}{W/m^2}$, about \qty{66}{mW/m^2}, some \qty{72}{\percent} of the
mean at the surface. (3) The heat crosses the boundary by conduction into
the lowest layer of the mantle; warmed, this rock becomes lighter than the
mantle above it and can rise, carrying the heat upwards: this is how plumes
are thought to start, deep down.
\end{solution}

\begin{solution}{pb:g12:internal-heat:1}
\textbf{1.} $4\pi (\num{6.371e6})^2 \approx \qty{5.101e14}{m^2}$, and
$2.073 + 3.028 = 5.101$: yes.

\textbf{2.} $3.028/5.101 \approx \qty{59}{\percent}$, about three fifths.

\textbf{3.} The ocean floor is young: it is made hot at the ridges and
cools as it moves away, and young \omterm{def:g11:lithosphere-asthenosphere:lithosphere}{lithosphere} loses much heat.

\textbf{4.} $490/\sqrt{4} = \qty{245}{mW/m^2}$.

\textbf{5.} $2.72 \times 0.026 \approx \qty{0.0707}{W/m^2} =
\qty{70.7}{mW/m^2}$: just below the mean, practically equal to it.

\textbf{6.} $0.0709 \times \num{2.073e14} \approx \qty{14.7}{TW}$.

\textbf{7.} $0.1054 \times \num{3.028e14} \approx \qty{31.9}{TW}$.

\textbf{8.} $14.7 + 31.9 \approx \qty{46.6}{TW}$, the \qty{47}{TW} of the
published estimate (the global mean of the chapter gives \qty{46.7}{TW}, the
value used below).

\textbf{9.} $\num{4.66e13}/\num{5.101e14} \approx \qty{0.0914}{W/m^2}$,
about \qty{91}{mW/m^2}.

\textbf{10.} $0.0243 \times \num{2.073e14} \approx \qty{5.0}{TW}$, a third of
the \qty{14.7}{TW} lost through the continents: the radioactivity of the
\omterm{def:g12:continental-crust:continental-crust}{continental crust} is a large part of their \omterm{def:g12:internal-heat:heat-flow}{heat flow}.

\textbf{11.} For each kilogram of uranium the Earth holds \qty{3.9}{kg} of
thorium, which release $3.9 \times 26.3 \approx \qty{102.6}{\micro W}$,
against \qty{94.4}{\micro W} for the uranium-238. Uranium-238 thus makes
$94.4/(94.4 + 102.6) \approx 0.479$ of the \qty{20.0}{TW}: about
\qty{9.6}{TW}. Uranium-235: $9.59 \times 4.10/94.4 \approx \qty{0.42}{TW}$.

\textbf{12.} $20.0 + 0.42 + 4 \approx \qty{24.4}{TW}$.

\textbf{13.} Granite and the other rocks of the \omterm{def:g12:continental-crust:continental-crust}{continental crust}, about
twenty times richer than the mantle's ultramafic rocks.

\textbf{14.} Larger: each isotope was more abundant then, twice as much for
every half-life back; 3 billion years ago the total was more than twice
today's.

\textbf{15.} $46.7 - 24.4 = \qty{22.3}{TW}$.

\textbf{16.} From the heat stored since the Earth formed: the \omterm{def:g12:internal-heat:secular-cooling}{secular
cooling} of the mantle and the core.

\textbf{17.} $22.3/46.7 \approx \qty{48}{\percent}$.

\textbf{18.} From $11/46.7 \approx \qty{24}{\percent}$ to $38/46.7 \approx
\qty{81}{\percent}$: the geoneutrinos narrow this range a great deal.

\textbf{19.} \omterm{def:g12:internal-heat:convection}{Mantle convection}, which carries most of the heat; only the last
tens of kilometres are crossed by conduction.

\textbf{20.} $24.4/46.7 \approx 0.52$: radioactivity supplies
\textbf{about half} of the Earth's heat loss today, and the slow cooling of
the planet the other half.
\end{solution}
